\section{Results I: the published results reproduce}
\label{sec:reproduction}

\paragraph{Per-task reproduction (RQ1).} With the benchmark flags (\texttt{tfm3-eval}), our SQL is within $\pm$3\% of Google's
published value on \AReproShare\% of the \ATasks{} tasks we could score (Table~\ref{tab:tiera}); the median deviation is zero to four
decimals and the deviations show no systematic sign. The two remaining tasks of the sample (\texttt{boomlet\_619},
\texttt{boomlet\_1209}) exhaust the 15\,GB of the T4 GPU under symmetric averaging, which doubles the batch, and are not scored;
the package defaults (\texttt{tfm3}) run on all 30. The five tasks outside the $\pm$3\% band (deviations from $-10.6\%$ to
$+6.1\%$; four of the five in our favour) all have covariates: the median absolute deviation is 0.001\% on tasks without covariates
and 1\% on tasks with covariates. We did not trace the cause; candidates are the handling of covariates with missing values and
batch-size choices that depend on the hardware. The per-task values are in Appendix~\ref{app:tiera}.

\begin{table}[!htbp]
\centering
\caption{Tier A: reproduction of Google's published fev-bench results. ``$\pm$3\%'': share of tasks whose SQL is within 3\%
of the published value; ``Median'' and ``Range'': deviations relative to the published SQL. ``Ratio'': geometric mean over tasks of SQL(TimesFM~3) /
SQL(Chronos-2), with our Chronos-2 (``ours'', no cross-learning) or fev's published Chronos-2 (``publ.'', cross-learning).}
\label{tab:tiera}
\scriptsize
\setlength{\tabcolsep}{3.5pt}
\begin{tabular}{llrrrlrr}
\toprule
Model & Tasks & $n$ & $\pm$3\% & Median & Range & Ratio ours & Ratio publ. \\
\midrule
TimesFM-3 eval & all tasks & 28 & 82\% & 0.00\% & [-10.6\%, +6.1\%] & 0.954 & 0.973 \\
TimesFM-3 eval & tasks passing both controls & 10 & 90\% & 0.00\% & [-0.4\%, +6.1\%] & 1.004 & 0.997 \\
TimesFM-3 & all tasks & 30 & 77\% & +0.24\% & [-8.9\%, +29.9\%] & 0.970 & 0.972 \\
TimesFM-3 & tasks passing both controls & 12 & 75\% & +0.55\% & [-0.3\%, +29.9\%] & 1.027 & 0.990 \\
\bottomrule
\end{tabular}

\end{table}

\paragraph{The claim against Chronos-2.} Google's comparison sets its TimesFM~3 results against the Chronos-2 results published by
fev. Our Chronos-2 failed the pipeline control on 16 of the 28 tasks (two further tasks failed the Seasonal Naive control), with a
higher SQL than published on 12 of them, so we traced the cause. fev's
published Chronos-2 results were produced with \emph{cross-learning} switched on: the series of a batch of up to 100 share
information through Chronos-2's group attention. Running fev's own evaluation script on five failing tasks reproduces the published
SQL to a relative difference of $4\cdot10^{-6}$ with cross-learning on and our values with it off (Appendix~\ref{app:chronos}); input preparation (batching,
imputation, timestamps) changes nothing. TimesFM~3 has no comparable option: it forecasts each series of a task on its own, or the
series of one multivariate task jointly. Cross-learning lowered Chronos-2's SQL by 2--9\% on four of the five tasks and raised it by
11\% on the fifth. With both models forecasting each series on its own, the geometric-mean SQL ratio of TimesFM~3 to Chronos-2 is
\ASkillOurs{} on all scored tasks, against \ASkillPubl{} in the published comparison: without cross-learning on the Chronos-2 side,
TimesFM~3's lead is slightly larger, not smaller. The lead comes from the \AFailTasks{} tasks where the controls fail
(\ASkillOursFail{} ours, \ASkillPublFail{} published); on the \AControlTasks{} tasks where our Chronos-2 matches fev's, the two
models are at parity in both comparisons (\ASkillOursCtl{} and \ASkillPublCtl{}). The published claim therefore holds on this
sample; what a leaderboard does not show is that the two result sets were produced with different inference options, one of which
favours Chronos-2 on most of the tasks it affects.

\paragraph{Against classical methods (H1, Tier A side).} With Google's inputs, TimesFM~3 (\texttt{tfm3}) has a geometric-mean
MASE of \HoneAMase{} and SQL of \HoneASql{} relative to fev's published statistical ensemble over the \HoneATasks{} tasks for which fev reports the ensemble, and a lower
SQL on \HoneAWins\% of them. Section~\ref{sec:postrelease} sets these numbers against the same comparison on post-release data.
